% A.4.57: explicit finite exclusion proof, pending independent final audit.
% The complete genuine-source shadow enumeration is a separate computation.
\subsection{The common finite CM comparison and its separating five-point level}

We specify the finite level before choosing its parameter.
Let $W$ be the actual ordered-pole quotient (FBO.13), and put
\[
                         s=12^3\,32\cdot3^{14}.
                                                               \tag{FCP.1}
\]
The regular representation of $W$ proves that $W_s$ in
(BCV.1) maps onto $W$.  Their ordinary outer images have
the actual quotient $p_W:J_s\to J_W$ induced by the
identity on $\Gamma_{1,2}$.
Use the integers and maps (BCV.1)--(BCV.6) for this $s$,
including $v,z,q_Z$ and $\eta_I:\mathcal M\to E_I$.
Write $\mathcal I_9(d)$ for the explicit bound (ICR.4),(ICR.7)
at $n=9$, starting at $d\,9!\,2^{100}$ to include the
label and relative-homology join.  Put
\[
                      e_I=\mathcal I_9(9!z^{10}).
                                                               \tag{FCP.2}
\]
This uses the seventeen specified recursion steps and gives
$|E_I|\leq e_I$.

At the actual configuration $E[3]$, integral linear maps
give the ordinary homomorphic section
\[
 s_{\mathrm{lin}}:\operatorname{SL}_2(\mathbb Z)\to\mathcal M,
 \quad A\mapsto(z\mapsto Az),\qquad
 \Omega_I=\eta_I\circ s_{\mathrm{lin}}.
                                                               \tag{FCP.3}
\]
They permute $E[3]$ and fix zero.  Their image is the
zero-fixing part $\mathcal N_0$ of the ordinary translation
normalizer; forgetting the other eight labels is its
inverse by (A.C4).
The period coordinate induced from $\Gamma(6)$ in A.4.56
has kernel $\ker\Omega_I\cap\Gamma(6)$ and hence maps
onto this one by the identity on $\operatorname{SL}_2(\mathbb Z)$:
both subgroups are normal, so their intersection is the
normal core in that induction.

Define the actual finite map
\[
 j_s=\rho_s\operatorname{mon}\delta s_{\mathrm{lin}}:
             \operatorname{SL}_2(\mathbb Z)\to J_s.
                                                               \tag{FCP.4}
\]
The central $-I$ is the elliptic involution and has trivial
branch action.  Thus $j_s$ is $\rho_s\varphi$ after
rigidification to $B=C_2*C_3$, through the ordinary geometric
map (A.C6)--(A.C7), with its actual marking.

Let $\mathcal M^0$ be the label-zero stabilizer, of index nine,
and let $\operatorname{forg}_0:\mathcal M^0\to\Gamma_{1,1}$
be ordinary forgetting.  Induce the actual homomorphism
$j_s\operatorname{forg}_0$ through those nine cosets.
Its identity coordinate on the stabilizer is precisely that
map, and its image has order at most
$9!|J_s|^9\leq9!w_s^{27}$.
Take the actual joint source image
\[
 \begin{split}
 q_R:\mathcal M&\longrightarrow Q_R,\\
 Q_R&=\operatorname{im}\bigl(q_{G^\dagger},\eta_I,
       \Omega_I\operatorname{per},
       \operatorname{Ind}_{\mathcal M^0}^{\mathcal M}
                    (j_s\operatorname{forg}_0)\bigr).
 \end{split}                                                     \tag{FCP.5}
\]
The first coordinate is the fixed actual map of (ECR.1).
If $D^\dagger=|G^\dagger|$, a specified finite integer from
that actual presentation, then
\[
                    |Q_R|\leq d_R
                      :=D^\dagger e_I^2\,9!w_s^{27}.
                                                               \tag{FCP.6}
\]
The induced coordinate uses the actual coset transitions;
its nine values have not been chosen independently.

Apply A.4.47 to $q_R$, and again to the resulting quotient,
adjoining (ECR.21).  This gives actual maps
\[
 \mathcal M\longrightarrow E_1^*\longrightarrow E_0
                            \longrightarrow Q_R.
                                                               \tag{FCP.7}
\]
Its proof applies to this refinement: the initial subgroup
is $\ker(q_R|_{K_0})$, the later subgroups are the same
characteristic cohomology-killing covers, and the retained
old factors supply the label-period comparison.
The additional factors are retained in that proof; no new
centralizer equality for them is assumed.
With the function $\mathcal R$ in (CFF.2), set
\[
                    b_*=24\mathcal R(\mathcal R(d_R))^{25}.
                                                               \tag{FCP.8}
\]
It bounds $|E_1^*|$.
Retain the original orientation input $s_0=128\cdot3^{1172}$
and the modular source detector of A.4.2, and define
\[
 \begin{gathered}
 m_\Delta=2^{5038380}3^{1371194},\qquad
 d_\Delta=(m_\Delta!)^{(m_\Delta!)^2},\\
 M=\max\{b_*,m_{v,9},m_{s_0,9},w_s^3,d_\Delta,2\},\\
 \mathcal Z=Z_1(M)
             =\Gamma_{0,5}/\kappa_{N(M)}(\Gamma_{0,5}).
 \end{gathered}                                                   \tag{FCP.9}
\]
Every function here is the integer expression already
displayed in (BCV.1)--(BCV.3), (ICR.4),(ICR.7),
(ECR.2)--(ECR.5), and (FGA.1)--(FGA.3).
Thus this specifies one characteristic finite quotient of
the ordinary five-point group and its full profinite
completion.  All source records can be joined with the
actual action on $B/\kappa_{m_\Delta}(B)$, where $\Delta$
is defined by A.4.2.

\begin{result}{A.4.57}{Proved}{An explicitly specified finite five-point level excludes every genuine parity-one pentagon shadow}
For the actual group $\mathcal Z$ in (FCP.9), every genuine
gentle element $F\in\widehat{\mathrm{GT}}_0$ satisfies
\[
            p(F)=1\text{ in }\mathcal Z
                       \quad\Longrightarrow\quad\Delta(F)=0.
                                                               \tag{FCP.10}
\]
The proof supplies one conjugator in the ordered-pole finite
target and its coupled signed modulo-four and modulo-three exponents, for
every multiplier.  This is an explicit finite exclusion
theorem.  The complete enumeration of the genuine source
records at this enormous level has not been executed and
is not asserted by this result.
\end{result}

\begin{proof}
\par\medskip\noindent\textbf{The exact multiplier-one transport domain.}\quad
First let $F=(1,f)$ and take the actual finite actions and
surface transports supplied by its pentagon in $\mathcal Z$.
The bound (FCP.9) retains (FCP.7), the surface input $Z$
of (BCV.3), and the orientation input of
(CFF.18)--(CFF.22).
Use the nine-point quilt and its core correction to write
\[
 \Psi_F(\iota)=c_F\iota c_F^{-1},\qquad
 \Psi_F^0=\operatorname{Ad}(c_F^{-1})\Psi_F,
                       \quad\Psi_F^0(\iota)=\iota
                                                               \tag{FCP.11}
\]
in $U_{1,[9]}(M)$.
All powers in $c_F$ have ordinary integral representatives
in this finite group.  Choose an actual pure core
multitwist lift $c\in\mathcal M$; every retained map factors
through the same finite input, so the lift has the required
value on all of them.

The proof of A.4.49 at the joined output (FCP.7) supplies
one ordinary pure positive $g\in\mathcal M$ such that
\[
       \Phi=\operatorname{Ad}(\bar g^{-1})\Psi_F
                         \quad\hbox{fixes }H\hbox{ in }E_0.
                                                               \tag{FCP.12}
\]
Its inputs are exactly the compatible finite automorphisms,
the peripheral power-one transport, the actual label-period
fibre, and the 289 specified covers.  The reduced cone
gives a fixed configuration; the four compact-matrix
obstruction rows and the separate peripheral determinant
force its grid sign positive.  The added factors only
refine the initial cover and are retained by that same
construction.  No automorphism of $E_0$ is assumed.

\par\medskip\noindent\textbf{One correction from translation fusion to the branch centralizer.}\quad
The ordinary involution normalizes $H$.  Thus $\Phi(\iota)$
normalizes $H$ in $E_0$, and recognition projects it to
the image of an ordinary $n\in N_{\mathcal M}(H)$ under
every factor of $q_R$.
Its label permutation is $h\mapsto-h$: the gentle action
preserves labels and $g$ is pure.  Hence $n$ fixes zero
and is $s_{\mathrm{lin}}(A)$ for an integral matrix $A$.
The compact forgetting action fixes $-I$, because the
genuine three-strand central word maps to its $\lambda$th
power and $\lambda=1$.  Quilt changes and conjugation by
the compact period of $g$ do not alter this central value.
Therefore $\Omega_I(A)=\Omega_I(-I)$, giving
\[
                         \eta_I(\Phi(\iota))=\eta_I(\iota).
                                                               \tag{FCP.13}
\]
With (FCP.11), the actual image of $g^{-1}c$ consequently
centralizes $\iota$ in $E_I$.
Equation (BCV.6) supplies one ordinary $d\in\mathcal C$
such that
\[
                         q_Z(d)=q_Z(g^{-1}c).
                                                               \tag{FCP.14}
\]
This is the one change from translation correction to branch
correction.  It is an equality in the retained ambient group,
not separate choices at two and three.

For the homomorphism $\mathcal B_F$ in (BCV.8), projection
of $\Phi=\operatorname{Ad}(g^{-1}c)\Psi_F^0$ now gives
\[
 \theta_s(q_Z(\Phi(a)))=D_s\mathcal B_F(a)D_s^{-1},
 \qquad D_s=\theta_s(q_Z(d))\in J_s\qquad(a\in\mathcal C).
                                                               \tag{FCP.15}
\]
Every term is in the stated domain: both $\Psi_F^0(a)$
and $\Phi(a)$ centralize the named involution in $E_I$
by (FCP.11),(FCP.13), and (BCV.6) projects them to
$q_Z(\mathcal C)$.
The branch-kernel proof (BCV.9)--(BCV.10) identifies
this homomorphism on every centralizer element with the
standard same-parameter automorphism $\beta_F$ of $J_s$:
\[
       \mathcal B_F(a)=
          \beta_F\bigl(\rho_s\operatorname{mon}\delta(a)\bigr).
                                                               \tag{FCP.16}
\]
The cone, peripheral-root and even-conjugator comparisons
are all used here.  Disk-generator surjectivity alone
would not give (FCP.16).

\par\medskip\noindent\textbf{Ordinary forgetting and one finite CM square.}\quad
For $a=s_{\mathrm{lin}}(A)\in\mathcal N_0$, its image
$\Phi(a)$ normalizes $H$ in $E_0$.
Recognition gives an ordinary zero-fixing normalizer
representative at $q_R$, because the label fixes zero.
For that ordinary representative the branch coordinate
and the induced forgetting coordinate are exactly the
same map $j_s$ in (FCP.4).  Their equality therefore
holds on this projected finite normalizer image.

Let $u_F\in J_s$ be the finite quilt weight
$T_F(Q_1,\operatorname{forg}_0Q_9)$ from the pinned
one-point torus quilt to the quilt obtained by forgetting
eight labels from the chosen nine-point quilt, evaluated
under $j_s$.  Substitution in (FGA.8)--(FGA.9) gives
$\Psi_{Q'}=\operatorname{Ad}(T_F(Q,Q')^{-1})\Psi_Q$.
Forgetting, followed by the identity induced coordinate
of (FCP.5), consequently gives
\[
 \begin{gathered}
 \theta_s(q_Z(\Phi(s_{\mathrm{lin}}(A))))
           =R_s j_s(\widetilde\alpha_F(A))R_s^{-1},\\
 R_s=j_s(\operatorname{forg}_0g)^{-1}u_F^{-1}.
 \end{gathered}                                                   \tag{FCP.17}
\]
Here $\widetilde\alpha_F$ is the genuine normalized
three-strand action on the full profinite $\Gamma_{1,1}$,
whose rigidification is the cusp-normalized gentle
automorphism $\alpha_F$ of $\widehat B$.
This is the pinned comparison supplied by A.4.40 through
the finite same-parameter relation.  Its single actual
quilt inner weight remains in (FCP.17).

Combine (FCP.15)--(FCP.17) and set $h_s=D_s^{-1}R_s$.
Since $j_s=\rho_s\varphi$ after rigidification and the
central involution dies under that map, the equality
factors through the actual source $\widehat B$:
\[
       \beta_F\rho_s\varphi
              =\operatorname{Ad}(h_s)\rho_s\varphi\alpha_F.
                                                               \tag{FCP.18}
\]
Projection to $J_W$ gives the finite CM comparison
\[
 \begin{gathered}
 \Phi_W=p_W\beta_F\rho_s:\widehat\Gamma_{1,2}\to J_W,
                         \qquad h_W=p_W(h_s),\\
 \Phi_W\varphi
               =\operatorname{Ad}(h_W)\rho_W\varphi\alpha_F.
 \end{gathered}                                                   \tag{FCP.19}
\]
The first map is an actual homomorphism with the standard
four-strand generator record of $F$.
No automorphism of $J_W$ is assumed: the automorphism
is $\beta_F$ on the retained input $J_s$, followed by
the specified output projection.

\par\medskip\noindent\textbf{The exact point-pushing quotient and its finite linearization.}\quad
The exponent lattice of $L_c$ in (QF.1) is $6\mathbb Z^2$.
Since $R_c=L_c\cap\Sigma'\Sigma^{12}$ and the second factor
is the full preimage of $12\mathbb Z^2$, the exponent lattice
of $R_c$ is exactly $12\mathbb Z^2$.
Thus $\mathcal P_c^{\mathrm{ab}}=(\mathbb Z/12)^2$, in its
displayed $a,b$ basis.  The actual joint quotient (FBO.13)
therefore has the explicit isomorphism
\[
 W\xrightarrow{\sim}\mathcal P_c\times\mathbb Z/12,
 \qquad A\mapsto(a,0),\quad B\mapsto(b,0),\quad C\mapsto(1,1).
                                                               \tag{FCP.L1}
\]
Indeed the first two abelian coordinates in (FBO.13) are
already the abelian image of its $\mathcal P_c$ coordinate;
the remaining coordinate is the exponent of $C$.
Conversely the three displayed generator images generate
this product, proving both directions of the isomorphism.

Every ordinary Artin action on this product has the form
\[
          (p,z)\longmapsto
                \bigl(\psi(p),z+\ell(p^{\mathrm{ab}})\bigr),
 \quad\psi\in\operatorname{Aut}(\mathcal P_c),\quad
 \ell\in\operatorname{Hom}((\mathbb Z/12)^2,\mathbb Z/12).
                                                               \tag{FCP.L2}
\]
This follows on the generators from (FBO.12)--(FBO.14),
and such maps are closed under composition and inverse.
The $\psi$ component is the actual folded modular action.
An inner automorphism of $W$ acts on the first factor by
an inner automorphism and is the identity on the central
second factor and on its exponent character.
Consequently the kernel of $J_W\to J_c$ is exactly the
row group of the maps (FCP.L2) with $\psi$ inner.
Removing that inner action gives the unique row $\ell$.
All rows occur: $\xi_1$ and $\xi_2$ have the primitive
rows $(0,-1)$ and $(1,0)$ in (FBO.14).
Thus there is the exact ordinary-image sequence
\[
 1\longrightarrow V=(\mathbb Z/12)^2
      \longrightarrow J_W\longrightarrow J_c\longrightarrow1,
 \qquad q_W(\xi_i)=e_i.
                                                               \tag{FCP.L3}
\]
The subgroup generated by the two modular braid generators
has row zero in (FCP.L2) and supplies a section $J_c\to J_W$.
Its kernel consists exactly of the inner $\psi$, which are
already inner on $W$.  Hence this is the actual split image
$J_W=V\rtimes J_c$; in particular $|J_W|=144|J_c|$.

The action on $V$ is the signed compact period, not an
independent module action.  A column vector $v$ in the
$\xi_1,\xi_2$ basis corresponds to the row
$v^{\mathsf T}J_0$, where
$J_0=\left(\begin{smallmatrix}0&-1\\1&0\end{smallmatrix}\right)$.
Conjugation by a period matrix $P$ sends that row to
$v^{\mathsf T}J_0P^{-1}=(Pv)^{\mathsf T}J_0$, since
$P^{\mathsf T}J_0P=J_0$ over $\mathbb Z/12$.
This proves the stated signed action through the actual
projection $J_c\to\operatorname{SL}_2(\mathbb Z/12)$.

Now use the standard finite braid record of $\Phi_W$.
At multiplier one it has
\[
 \Phi_W(a)=a,\quad
 \Phi_W(b)=f(a^2,b^2)^{-1}bf(a^2,b^2),\quad
 \Phi_W(c)=e^{-1}ce,\qquad e=f(ba^2b,c^2).
                                                               \tag{FCP.L4}
\]
Every term here is evaluated through the specified quotient.
Before that evaluation the proword $e$ lies in the actual
completed point-pushing kernel $K_{\mathrm{push}}$:
under folding, $c\mapsto a$ and
$ba^2b=(ab)^3a^{-2}$ commutes with $a^2$, so its two
arguments commute and the commutator proword is one.
Moreover $a$ commutes with both arguments of $e$ in
$\widehat\Gamma_{1,2}$, so $[e,a]=1$.
Thus its finite image is a vector $E\in V$ with
$(U-I)E=0$.
Folding the three displayed images gives the genuine
three-strand record.  In particular the two point-pushing
basis words have images in $\ker(J_W\to J_c)=V$.
Since $\Phi_W$ is a homomorphism and $V$ is abelian of
exponent twelve, its restriction factors through
$K_{\mathrm{push}}^{\mathrm{ab}}/12=V$ and gives an
actual endomorphism $N:V\to V$.

The first basis word gives
$Ne_1=e_1+(U-I)E=e_1$.
Covariance with $a$ gives $NU=UN$.
Writing a two-by-two matrix over $\mathbb Z/12$ shows
directly that
\[
             N=\begin{pmatrix}1&t\\0&1\end{pmatrix}=U^t
                         \quad(t\in\mathbb Z/12).
                                                               \tag{FCP.L5}
\]
In particular it is invertible; that property has not
been inferred from an unproved automorphism of $J_W$.
Covariance with $b$ now yields
\[
        NLN^{-1}=q_{12}(f)^{-1}Lq_{12}(f).
                                                               \tag{FCP.L6}
\]
The R1 signed covariance gives the same equation for
$V_F=U^{8\rho}$, with both $N$ and $V_F$ fixing $e_1$
and commuting with $U$.
Their quotient centralizes $U,L$; entrywise multiplication
over $\mathbb Z/12$ makes it scalar, and its value on
$e_1$ makes that scalar one.
Therefore $N=U^{8\rho}$, the precise finite version
of (A.C9).  This proves the required linearization
through its actual homomorphism and kernel.

\par\medskip\noindent\textbf{The finite signed CM calculation.}\quad
Use the genuine pure frames $\beta,A$ and shear
$V=U^{8\rho}$ of (A.C1)--(A.C3).
Evaluate the same word $A$ in the two-point braid
generators and conjugate the left side of (FCP.19)
by that value.  Its target conjugator is now the
single element $h=\widetilde A^{-1}h_W\in J_W$.

The calculations (A.C8)--(A.C13) apply to these specified
finite outputs.  The braid words are evaluated under the
actual homomorphism $\Phi_W$, and point-pushing homology
is retained with coefficient one by (FBO.14).
Ordinary forgetting retains the whole actual $J_c$
factor.  Pure frame covariance and the signed
abelianization modulo twelve make both sides induce
the identity on the same cyclic CM subgroup of $J_c$.
Thus the $J_c$ image of $h$ centralizes its named CM
generator.  The finite centralizer calculation A.4.18
and the ordinary outer-image identification A.4.22
give two residues $i,j\in\mathbb Z/6$ attached to this same
element, with the exact joint constraint
\[
 \operatorname{per}_4(h)=R_+^i,\qquad
 \operatorname{per}_3(h)=R_+^j,\qquad i\equiv j\pmod2.
                                                               \tag{FCP.20}
\]
These are the eighteen actual residues in (QH.4)--(QH.5).
No stronger claim of a single cyclic modulo-twelve value
is needed or made.
The remaining kernel of this finite period map acts
trivially on point-pushing homology, by the literal
conjugation table (A.C8),(FBO.14).

The ordinary map $\varphi$ supplies its integral
commutator-homology marking $C$, with
$CR=R_+^\epsilon C$ and the actual sign
$\epsilon\in\{1,-1\}$ retained in (A.C7).
Applying point-pushing homology modulo twelve to
(FCP.19), after the frame correction, gives
\[
               TC=R_+^iCM\pmod4,\qquad
               TC=R_+^jCM\pmod3,
                                                               \tag{FCP.21}
\]
where $M$ is the genuine source commutator-homology
matrix (A.C11)--(A.C13).
At three its scalar on the specified CM line is one,
so the second equation gives $\epsilon_3(T)=j\pmod2$.
At four, $T=I$ and the source matrix is exactly
$(-1)^{\epsilon_2(W)}R^\ell$ with $\ell$ even; the same
first equation gives $\epsilon_2(W)=i\pmod2$.
The sign $\epsilon$ changes neither parity.  Hence
$\Delta(F)=\epsilon_3(T)+\epsilon_2(W)=i+j=0$ in
$\mathbb F_2$, by the actual joint constraint (FCP.20).
The executed certificate
\path{check_finite_cm_parity_bridge.py}
checks these finite equations for all eighteen coupled
residues; it is the paired-exponent replacement for the
earlier single-exponent matrix check.
Its use here follows the actual square (FCP.19).

\par\medskip\noindent\textbf{All multipliers through actual record composition.}\quad
Take any genuine $F$ whose pentagon is one in $\mathcal Z$.
Cyclotomic surjectivity gives an actual
$\sigma\in G_{\mathbb Q}$ with
$\lambda(\sigma)=\lambda(F)$.
The Galois embedding is into the full profinite
$\widehat{\mathrm{GT}}$, and the refereed character
satisfies $\Delta(\sigma)=0$.
The actual gentle element $F_0=F\sigma^{-1}$ has
multiplier one and $\Delta(F_0)=\Delta(F)$.

On every universal finite tower group use the actual
composed automorphism $\Psi_F\Psi_\sigma^{-1}$.
Compose the compatible surface transports and naturality
comparisons in the same order.  Both factors exist at
the common level: $F$ by its assumed pentagon in
$\mathcal Z$, and $\sigma$ by its full pentagon.
The literal finite record law (FGA.10)--(FGA.12),
including the coefficient cocycle, identifies every
local path weight and braid-generator record of this
composite with those of $F_0$.
Its peripheral powers are exactly one.

The preceding proof uses exactly those compatible finite
automorphisms, the actual peripheral surface covariance,
and the local generator weights.
The translation-fusion proof used in (FCP.12) uses these
same action properties and the fixed finite cohomology
covers; its orientation check now has multiplier one.
Thus (FCP.11)--(FCP.21) apply to this actual composed
transport and give $\Delta(F_0)=0$, hence $\Delta(F)=0$.
There is no assertion that $F_0$ belongs to the original
reduced pentagon locus, or that that locus is a group.
The actual composition supplies precisely the action
data used in the proof.
\end{proof}

The identity parameter has parity zero and pentagon one
at this level, so the set of parities of its genuine
pentagon shadows is exactly $\{0\}$.
The theorem proves finite exclusion through an explicit
source construction and exact finite comparisons.
Generating every genuine finite source record, listing
all its pentagon values and attaching its parity still
requires an executed enumeration.  The explicit size
bounds and finite matrix checks have not been presented
as that enumeration.
